The results support three main observations.

First, open-set RF preliminary-action recognition should be evaluated as an operating-point problem, not only as a ranking problem. Table~\ref{tab:main_osr_results} shows that DQNGuard improves unknown F1 and OSR macro F1 while maintaining low known rejection under the fixed surrogate-open condition. This distinction matters because an OSR layer can appear effective if it rejects many samples, but such behavior would reduce the value of the known PA evidence stream. DQNGuard's known-budgeted thresholding makes the comparison stricter: unknown detection improves while the allowed loss of known PA utility remains constrained.

Second, surrogate-open calibration is useful but target-dependent. Figure~\ref{fig:target_surrogate_matrix} shows that the surrogate class is not an interchangeable source of generic ``unknownness'' evidence. Some surrogate behaviors induce guard evidence that transfers well to a target unknown, while others provide little useful separation. This suggests that surrogate-open calibration should be treated as a deployment assumption to be tested, not as a universal substitute for target-open calibration.

Third, the RF front end should be interpreted as a sensing and triage component. DQNGuard does not infer a complete attack chain, select an EW response, or assign final semantic labels to all novel emissions. Its narrower role is to decide whether an OTA RF window should be emitted as known PA evidence or preserved as a nonconforming behavior for later analysis. That boundary is important: the value of the method is not that it solves the downstream QR-CWoS policy problem, but that it supplies a cleaner evidence interface to such systems.

The main limitation is surrogate selection. The Target--Surrogate Matrix demonstrates that some surrogates transfer better than others, but the current experiments do not provide a reliable automatic rule for choosing the best surrogate before the target unknown is observed. A practical deployment would therefore require either a conservative fixed-surrogate policy, pooled surrogate evidence, human-reviewed surrogate diagnostics, or a learned surrogate-selection rule trained across a broader behavior taxonomy.

A second limitation is scale. The evaluated corpus spans three protocol families and five PA behaviors, which is sufficient to expose target-dependent surrogate transfer but not sufficient to characterize the full space of RF behavior novelty. Larger PA taxonomies would allow leave-multiple-out evaluation, pooled calibration studies, and stronger tests of whether the observed transfer asymmetries are specific to the current behaviors or reflect a broader property of RF open-set recognition.

Finally, the evaluation begins after RF window extraction and feature construction. Operational performance will also depend on upstream burst detection, segmentation, synchronization, and OTA capture quality, as well as downstream label-making and response-policy modules. These dependencies define the boundary of the claim. The result shown here is that unknown PA behaviors can be detected from OTA RF windows under controlled calibration regimes, while preserving a bounded stream of known PA evidence.
